% IOE BE Project LaTeX Template
% Final Report — adapted from the Mid-Term Defense template
% (Documentation/97-Educational_Contents/EdumentX_Mid_Defense/).
% New chapters 4-7 added for features implemented after mid-term.

\documentclass[12pt]{report}
\usepackage{xcolor}
\usepackage[utf8]{inputenc}
\usepackage[numbers,sort&compress]{natbib}         % numeric [1] citations
\usepackage[a4paper,left=1.4in,right=1.0in,top=1in,bottom=1in]{geometry}
\usepackage{float}                                 % [H] figure placement

% Digital-copy hyperref (color links so nav panel is usable):
\usepackage[hidelinks]{hyperref}

\linespread{1.3}
\usepackage{graphicx}                             % figures
\usepackage{rotating}                             % landscape figures
\usepackage{amsmath}                              % math
\usepackage{amssymb}                              % extra math symbols
\usepackage{titlesec}                             % chapter/section formatting
\usepackage{fontspec}                             % XeLaTeX/LuaLaTeX font selection
\IfFontExistsTF{Times New Roman}{%
  \setmainfont{Times New Roman}%
}{%
  \setmainfont{Liberation Serif}%
}
\usepackage{enumitem}                             % nicer itemize spacing
\usepackage{longtable}                            % multi-page tables
\usepackage{booktabs}                             % better-looking table rules
\usepackage{microtype}                            % better hbox/line breaking
\usepackage{ragged2e}                             % \RaggedRight to fight overfull hbox

% --- Spacing / line breaking ----
\setlength{\parskip}{0.5\baselineskip plus 1pt minus 1pt}
\setlength{\parindent}{0pt}
\setlength{\emergencystretch}{2em}
\sloppy
\tolerance=2000
\hbadness=10000
\hfuzz=2pt

% --- Chapter / section spacing -----
\titlespacing*{\chapter}{0pt}{1.0\baselineskip}{0.8\baselineskip}
\titlespacing*{\section}{0pt}{0.7\baselineskip}{0.3\baselineskip}
\titlespacing*{\subsection}{0pt}{0.5\baselineskip}{0.2\baselineskip}
\titleformat{\chapter}[hang]
  {\normalfont\Large\bfseries}
  {\chaptertitlename\ \thechapter.}{1em}{}
\renewcommand{\chaptername}{}

% --- Itemize spacing ---
\setlist[itemize]{leftmargin=1.2em,itemsep=0.15\baselineskip,topsep=0.2\baselineskip,parsep=0pt}

% --- Image search path ---
\graphicspath{{Images/}}

% --- safeincludegraphics wrapper ---
\providecommand{\safeincludegraphics}[2][]{%
  \includegraphics[#1]{#2}%
}

% --- \outputscreen{file}{caption}{label} ---
\providecommand{\outputscreen}[3]{%
  \begin{figure}[!htbp]
    \centering
    \includegraphics[width=0.85\textwidth]{#1}%
    \caption{#2}\label{#3}%
  \end{figure}%
}

\begin{document}
\begin{titlepage}
    \begin{center}
        \vspace*{0.2cm}
        {\includegraphics[scale=1.4]{image.png}}\\[0.8cm]
        {\large \MakeUppercase{Tribhuvan University}}\\[0.15cm]
        {\large \MakeUppercase{Institute of Engineering}}\\[0.15cm]
        {\large \MakeUppercase{National College of Engineering}}\\[1.0cm]

        {\normalsize A}\\[0.15cm]
        {\normalsize MINOR PROJECT REPORT}\\[0.15cm]
        {\normalsize ON}\\[0.6cm]
        {\large \MakeUppercase{EdumentX: Location Based Tutor Discovery Platform}}\\[1.2cm]

        {\textbf{SUBMITTED BY:}}\\[0.3cm]
        {\normalsize APOORVA PAUDEL (NCE080BCT006)}\\[0.12cm]
        {\normalsize ASIM POUDEL (NCE080BCT008)}\\[0.12cm]
        {\normalsize DIBYA SHAKYA (NCE080BCT013)}\\[0.12cm]
        {\normalsize SUHAN KHADKA (NCE080BCT043)}\\[1.0cm]

        {\textbf{SUBMITTED TO:}}\\[0.3cm]
        {\normalsize DEPARTMENT OF ELECTRONICS \& COMPUTER ENGINEERING}\\[2.5cm]

        {\normalsize August, 2026}
    \end{center}
\end{titlepage}


% --- Front matter --------------------------------------------------
\pagenumbering{roman}
\setcounter{page}{2}

% --- Page of Approval ---
\chapter*{Page of Approval}
\addcontentsline{toc}{chapter}{\numberline{}Page of Approval}
\vspace{1.5cm}
\begin{center}
\uppercase{
Tribhuvan University\\
Institute of Engineering\\
National College of Engineering\\
Department of Electronics and Computer Engineering
}
\end{center}
\vspace{.5cm}
The undersigned certifies that they have read and recommended to the Institute of Engineering for acceptance of a project report entitled \textbf{``EdumentX: Location Based Tutor Discovery Platform"} submitted by \textbf{Apoorva Paudel}, \textbf{Asim Poudel}, \textbf{Dibya Shakya}, \textbf{Suhan Khadka} in partial fulfillment of the requirements for the Bachelor's degree in Electronics, Information and Communication Engineering \& Computer Engineering.\\

\vspace{1cm}
\begin{minipage}{.5\textwidth}
\raggedright
.............................\\
Supervisor\\
\textbf{Subash Pandey}\\
Designation\\
Department of Electronics and Computer Engineering,\\
National College of Engineering, IOE, TU.\\
\end{minipage}%
\begin{minipage}{0.5\textwidth}
\raggedleft
.............................\\
External examiner\\
\textbf{Full Name}\\
Designation \\
Affiliation\\
\end{minipage}\\
\vspace{1cm}
\centerline{Date of approval:}

% --- Copyright ---
\chapter*{Copyright}
\addcontentsline{toc}{chapter}{\numberline{}Copyright}
The author has agreed that the Library, Department of Electronics and Computer Engineering, National College of Engineering, and Institute of Engineering may make this report freely available for inspection. Moreover, the author has agreed that permission for extensive copying of this project report for scholarly purposes may be granted by the supervisors who supervised the work recorded herein or, in their absence, by the Head of the Department wherein the project report was done. It is understood that recognition will be given to the author of this report and the Department of Electronics and Computer Engineering, National College of Engineering, IOE for any use of the material of this project report. Copying publication or the other use of this report for financial gain without the approval of to the Department of Electronics and Computer Engineering, National College of Engineering, Institute of Engineering, and the author's written permission is prohibited.\\
Request for permission to copy or to make any other use of the material in this report in whole or in part should be addressed to:\\\vspace{1cm}
Head\\
Department of Electronics and Computer Engineering\\
National College of Engineering\\
Talchhikhel, Lalitpur.

% --- Acknowledgments ---
\chapter*{Acknowledgments}
\addcontentsline{toc}{chapter}{\numberline{}Acknowledgements}
We would like to express our sincere gratitude to our project supervisor \textbf{Subash Pandey} for their invaluable guidance, continuous support, and constructive feedback throughout the development of this project. Their expertise and encouragement were instrumental in shaping the direction of EdumentX.

We are also grateful to the Head of the Department of Electronics and Computer Engineering, \textbf{National College of Engineering}, and all the faculty members who provided us with the knowledge and skills necessary to undertake this project.

We extend our thanks to our teammates \textbf{Apoorva Paudel}, \textbf{Asim Poudel}, and \textbf{Dibya Shakya} for their collaboration, dedication, and mutual support during every phase of this project. Finally, we thank our families and friends for their encouragement and patience throughout this journey.

% --- Abstract ---
\chapter*{Abstract}
\addcontentsline{toc}{chapter}{\numberline{}Abstract}
Finding a qualified home tutor in the Kathmandu Valley remains a fragmented, broker-dependent process. Students and parents rely on word-of-mouth referrals, social media groups, or local advertisements, resulting in mismatched expectations, wasted time, and limited trust in tutor qualifications. Existing tutoring platforms are either online-class-focused or lack location-based discovery, leaving a gap for in-person home tutoring services. EdumentX is a mobile application built with React Native and Expo that connects students and parents with verified home tutors in their locality. The platform provides map-based tutor discovery using Google Maps, a distance-based ranking system that promotes tutors by proximity and subscription tier, and a Retrieval-Augmented Generation (RAG) chatbot powered by Groq that answers natural-language queries about available tutors. Tutors undergo an AI-assisted identity verification pipeline that extracts text from citizenship documents using Groq Vision OCR and validates profile photographs, with automatic approval for high-confidence matches and manual admin review for ambiguous cases. The system supports enrollment management, group batch scheduling, in-app messaging, ratings and reviews, and an eSewa-based Pro subscription model for tutors. Firebase Authentication and Cloud Firestore provide identity and data persistence, while Supabase handles object storage and serverless edge functions. Security rules enforce owner-only data access and role-based permissions. The application runs on Android with Google Maps integration and is designed to operate within free-tier quotas, making it viable as a zero-cost demonstration platform. Testing includes 47 unit tests for enrollment logic, 21 lint rule tests for design token enforcement, and six Firestore security rules emulator suites. The results demonstrate that a location-first, trust-oriented tutoring marketplace can be built and deployed without paid infrastructure at demo scale.

\vspace{0.5cm}
\noindent\textit{Keywords: Location-Based Services, Tutor Discovery, React Native, Firebase, KNN Recommendation, RAG Chatbot, AI Verification}

% --- Table of Contents ---
\tableofcontents
\addcontentsline{toc}{chapter}{\numberline{}List of Figures}
\listoffigures
\addcontentsline{toc}{chapter}{\numberline{}List of Tables}
\listoftables
\chapter*{List of Abbreviations}
\addcontentsline{toc}{chapter}{\numberline{}List of Abbreviations}

\begin{tabular}{l l}
AI & Artificial Intelligence \\
API & Application Programming Interface \\
DFD & Data Flow Diagram \\
GPS & Global Positioning System \\
HMAC & Hash-based Message Authentication Code \\
HTTP(S) & Hypertext Transfer Protocol (Secure) \\
IOE & Institute of Engineering \\
JSON & JavaScript Object Notation \\
KNN & K-Nearest Neighbors \\
LBS & Location-Based Service \\
LLM & Large Language Model \\
NCE & National College of Engineering \\
NPR & Nepalese Rupee \\
OCR & Optical Character Recognition \\
RAG & Retrieval-Augmented Generation \\
REST & Representational State Transfer \\
SDLC & Software Development Life Cycle \\
SEE & Secondary Education Examination (Grade 10, Nepal) \\
TU & Tribhuvan University \\
UI & User Interface \\
UX & User Experience \\
\end{tabular}

\newpage

% ====================================================================
% CHAPTER 1: INTRODUCTION
% ====================================================================
\chapter{Introduction}
\pagenumbering{arabic}

\section{Background}

Education is one of the most important foundations for personal
growth and social progress. Around the world, families invest
heavily in their children's education, recognizing that formal
schooling alone may not always be enough. When students struggle
with certain subjects or need extra preparation for exams, many
families turn to private tutoring as a way to fill the gaps.
This practice of hiring tutors outside of regular school hours
has become extremely common, as parents look for additional
academic support to help their children perform better.

Private tutoring, commonly referred to as shadow education, has
become a widespread practice as formal schooling systems often
fail to meet the diverse learning needs of students
\cite{Bray2007Shadow}. This form of supplementary education
has grown particularly fast in developing countries, where
public school resources are limited and teacher-student ratios
remain inadequate \cite{Bray2012Shadow}. Research on shadow
education as a global phenomenon confirms that it has expanded
across all continents and income levels, becoming a parallel
curriculum that operates alongside formal schooling
\cite{Kim2019Shadow}.

Nepal's education system has made notable progress over the past
decade. According to the National Population and Housing Census
2021, the overall literacy rate reached 76.3 percent among the
population aged five years and above, up from 65.9 percent in
2011. The country has also experienced rapid urbanization during
this period, with the urban population share jumping from 17.07
percent in 2011 to 66.17 percent in 2021 \cite{CBS2023}.
Despite these improvements, challenges within the mainstream
education system persist. Overcrowded classrooms, inconsistent
teaching quality, and limited resources in public schools
continue to push families toward private tutoring. In cities
like Kathmandu, dual-income households have become increasingly
common, leaving working parents with limited time for direct
academic supervision of their children. As a result, many urban
families look for home tutors as a practical solution.

The phenomenon of shadow education is deeply rooted in Nepal's
educational landscape. Thapa's study on private and public
school competition in Nepal highlights how supplementary
tutoring has become an established practice, with families
across income levels relying on it to fill gaps left by formal
schooling \cite{Thapa2013Private}. Studies further indicate
that approximately 68 percent of Grade 10 students in Nepal
receive some form of private tutoring, confirming how
widespread this practice has become at the secondary level.

Despite this growing demand, Nepal's tutoring market remains
fragmented and informal. Parents and students typically rely on
word-of-mouth recommendations or unstructured social media
groups to find tutors. These methods offer no verification of
tutor credentials, no proximity-based filtering, and no
intelligent matching. With smartphone penetration in Nepal
reaching 65 to 71 percent, the digital infrastructure exists to
address these problems through a dedicated mobile platform.

This report presents EdumentX, a location-based mobile
application designed to modernize home tutor discovery in
Kathmandu Valley. The platform enables students and parents to
find nearby verified tutors through an interactive map, enroll
in long-term monthly plans, and receive personalized
recommendations through an AI-powered RAG chatbot.

\section{Problem Statement}

The process of finding a home tutor in Nepal is currently slow,
unreliable, and frustrating for everyone involved. Both students
seeking academic help and tutors looking for teaching
opportunities face significant difficulties due to the lack of
any organized system.

The core problems are:

{\renewcommand{\labelenumi}{(\roman{enumi})}
\begin{enumerate}
\item There is no way for parents to check whether a tutor is
actually qualified or trustworthy before hiring them. Most
families rely on recommendations from people they know, which
limits their options and gives no real measure of teaching
quality or safety.

\item Independent tutors, including university students looking
for part-time work, retired teachers, and subject specialists,
have no proper platform to advertise their services. Many end
up going through agencies that take large commissions, while
others simply cannot find students outside their own
neighborhood.

\item The tools people currently use to find tutors, mainly
WhatsApp groups and Facebook pages, are not designed for this
purpose. Posts are unorganized, location details are vague
(just area names with no actual distance), there is no way to
filter or search, and there is no system to verify anyone's
identity or qualifications.
\end{enumerate}
}

At present, no mobile application exists in Nepal that lets
users find tutors on a map, filter by distance, or get
AI-based suggestions. EdumentX aims to solve these problems
by offering proximity-based search on an interactive map, a
verification system where admins review tutor documents and
award Blue Tick badges, a structured monthly enrollment
process, and a chatbot that pulls live tutor data to give
personalized recommendations.

\section{Objectives}

\begin{itemize}
\item To develop a location-based mobile application that enables
students and parents to discover and connect with nearby home
tutors in Kathmandu Valley using map-based proximity search.

\item To integrate an AI-powered RAG chatbot that allows users to
find suitable tutors through natural language queries grounded
in live platform data.

\item To implement a tutor verification system that helps parents
make safe and informed hiring decisions.

\item To implement structured enrollment management, group batch
scheduling, in-app messaging, and a Pro subscription model for
tutors.
\end{itemize}

\section{Scope}

EdumentX is a role-based mobile platform for connecting students
and parents with home tutors based on location, subject, and
monthly enrollment preferences. The scope of this project
includes the following:

\begin{itemize}
\item \textbf{Target Users:} Students, parents, independent home
tutors, and platform administrators.

\item \textbf{Geographic Focus:} Kathmandu Valley, where
smartphone penetration and internet connectivity are
sufficient to support the platform.

\item \textbf{Service Type:} Physical (in-person) home tuition
services only; no online or video-based classes.

\item \textbf{Academic Levels:} School and college students at
SEE, +2, and undergraduate levels.

\item \textbf{Platform:} Android application built with React
Native (Expo SDK 54) and NativeWind 4.2.

\item \textbf{Discovery Methods:} Map-based search with proximity
ranking and an AI chatbot for natural language queries.

\item \textbf{Enrollment Model:} Long-term monthly enrollment
plans managed through structured request and response flows,
with slot-based scheduling.

\item \textbf{Verification:} AI-assisted identity verification
with automatic approval for high-confidence matches and manual
admin review for ambiguous cases. Blue Tick badge system.

\item \textbf{Group Learning:} Support for group tuition batches
created by tutors for shared, cost-effective learning.

\item \textbf{Subscription:} Pro Tutor subscription via eSewa
payment gateway (sandbox) for enhanced visibility and features.
\end{itemize}

% ====================================================================
% CHAPTER 2: LITERATURE REVIEW
% ====================================================================
\chapter{Literature Review}

\section{Related Work}

This section reviews recent academic research on location-based
tutor matching systems, AI-powered educational platforms, and
identifies the gaps that EdumentX aims to address.

A 2022 study presented at ASCILITE introduced MENTOR, an
intelligent mobile peer tutoring application designed for both
face-to-face and remote tutoring \cite{Chung2022MENTOR}. The
system focused on matching peer tutors with students based on
subject expertise and availability. While it demonstrated the
value of mobile platforms for tutoring, it did not include any
location-based features such as map visualization or proximity
ranking. The matching was based solely on academic parameters
without considering geographic distance between users.

A 2026 study published in the Journal of Emerging Technologies
and Innovative Research described a Flutter-based mobile
application for finding local tutors in India
\cite{Sharma2026LocalTutor}. The system used Google Maps API
for proximity-based tutor search and Firebase as the backend
database. Students could filter tutors by subject, availability,
and location. The application calculated distances between
users using the Haversine formula and included features for
booking and rating tutors. However, the researchers noted that
AI-based recommendations were left as future scope. The entire
matching process relied on manual filtering by the user, with
no intelligent ranking or conversational search capability.

In 2025, researchers from Malaysia developed iTeachU, a system
that combined AI with geolocation for tutor matching
\cite{Shari2025iTeachU}. The platform used content-based
filtering along with the K-Nearest Neighbors algorithm,
incorporating Haversine distance calculations to rank tutors
by both preference match and physical proximity. It was built
using Flutter for the frontend and Flask for the backend. While
the system demonstrated effective ranking within the Melaka
region, its AI component was entirely rule-based. Users could
not interact with the system using natural language. Instead,
they had to select from predefined filters and dropdown menus,
making the experience less flexible than a conversational
interface.

All of these studies represent important steps toward
technology-driven tutor discovery, but several gaps remain
when considering the Nepali context. No existing research or
platform provides geospatial visualization of tutors on
interactive maps with service area coverage for Nepal.
City-level dropdown menus, which are used by most platforms,
are not useful in dense urban areas like Kathmandu where a
tutor two kilometers away is very different from one ten
kilometers away. Independent part-time tutors such as
university students and retired teachers lack a proper digital
channel to reach nearby students, as existing systems tend to
favor agency-managed tuition centers. None of the reviewed
systems offer a natural language interface where users can
simply type what they need instead of navigating through
multiple filter options. While some platforms claim AI-based
matching, none combine proximity, qualifications, availability,
and budget into a single intelligent ranking powered by a
large language model. The current methods available in Nepal,
primarily social media groups, require significant manual
effort to browse posts, contact agencies, and compare options
without any automated assistance.

EdumentX brings map-based tutor visualization, a
conversational interface for natural language queries, a
structured verification system with Blue Tick badges, and a
focus on empowering independent tutors to connect with nearby
students efficiently.

\section{Related Theory}

\subsection{Location-Based Services}

Location-Based Services deliver context-aware services using a
device's real-time geographic position obtained through GPS,
cellular triangulation, or Wi-Fi. In EdumentX, location-based
services power the proximity-based tutor discovery and map
visualization features. Tutor coordinates are stored as
Firestore GeoPoint fields, which allow geospatial queries to
be executed directly from the client application.

\subsection{Haversine Distance Formula}

The Haversine formula computes the great-circle distance
between two points on Earth's surface, accounting for the
planet's curvature. EdumentX uses this formula to rank tutors
by their true distance from the student. The computation
proceeds as follows:

\begin{equation}
a = \sin^2\!\left(\frac{\Delta\phi}{2}\right) +
\cos(\phi_1)\cdot\cos(\phi_2)\cdot\sin^2\!\left(\frac{\Delta\lambda}{2}\right)
\end{equation}
\begin{equation}
c = 2\,\arctan2\!\left(\sqrt{a}\,,\,\sqrt{1-a}\right), \qquad
d = R \cdot c
\end{equation}

\noindent where $\phi_1, \phi_2$ are latitudes in radians,
$\lambda_1, \lambda_2$ are longitudes in radians and
$\Delta\phi$ and $\Delta\lambda$ are the coordinate
differences, $R \approx 6{,}371\;\text{km}$, and $d$ is the
output distance. Degree values from GeoPoints are converted
using $\phi_{\text{rad}} = \phi_{\text{deg}} \times
\frac{\pi}{180}$. For the 1 to 20 km range typical in
Kathmandu Valley, this formula provides meter-level accuracy.

\subsection{Geofencing and Radius-Based Filtering}

Before applying the Haversine formula, EdumentX uses a
rectangular bounding box pre-filter to narrow down the
candidate set and reduce computation
\cite{Shevchenko2023Geofencing}. For a given search radius
$r$ and student position $(\phi_s, \lambda_s)$:

\begin{equation}
\Delta\phi_{\max} = \frac{r}{R} \cdot \frac{180}{\pi}, \qquad
\Delta\lambda_{\max} = \frac{r}{R \cdot \cos(\phi_s)} \cdot \frac{180}{\pi}
\end{equation}

\noindent Only tutors whose coordinates fall inside this
bounding box proceed to the exact Haversine distance check,
which keeps query response times low.

\subsection{K-Nearest Neighbors Algorithm}

KNN is a non-parametric algorithm used in EdumentX to rank
tutors by proximity. The primary ranking criterion is
Haversine distance, with Pro subscription tier as a secondary
sort key that promotes subscribers above non-subscribers
within the same distance band. The default search radius is
15 km, which covers the Kathmandu Valley extent.

\subsection{Retrieval-Augmented Generation}

Retrieval-Augmented Generation is an approach that combines
information retrieval with text generation to produce responses
grounded in actual data rather than relying solely on a
language model's training data. The original RAG framework was
introduced by Lewis et al., who showed that combining a
pre-trained retriever with a sequence-to-sequence generator
significantly improves accuracy on knowledge intensive tasks
\cite{Lewis2021RAG}. In EdumentX, RAG is used to retrieve
matching tutor profiles from Firestore based on the user's
query and inject them as context into the Groq API (Llama 3
model) prompt, ensuring that recommendations are based on
real, live platform data. A recent survey on RAG chatbots in
education confirms that this approach reduces hallucination and
improves the relevance of generated responses in educational
contexts \cite{Swacha2025RAG}.

\subsection{React Native and Expo}

The EdumentX mobile app is built with React Native, which
lets a single codebase produce native iOS and Android
binaries. Expo sits on top of React Native and provides a
managed toolchain for device features such as location,
camera, and push notifications, along with over-the-air
updates that shorten the release cycle.

\subsection{Firebase and Supabase}

Four services form the EdumentX backend. Firebase
Authentication handles email/password and Google sign-in.
Cloud Firestore is the main NoSQL store; tutor locations are
saved as GeoPoint fields, and Security Rules control who can
read or write which documents. Supabase Storage (1 GB free)
holds verification documents and profile photos. Supabase
Edge Functions provide serverless compute for the RAG chatbot
(Groq API), eSewa payment signing, and AI identity
verification. The Firebase Admin SDK is used out-of-band from
a local Node.js script to grant the administrator role to a
fixed list of emails.

\subsection{Google Maps via expo-maps}

EdumentX uses Google Maps rendered natively through the
expo-maps module on Android, using the device's native Google
Play Services. Custom teardrop pins with tutor avatars are
rasterized offscreen and loaded into native map markers. The
Google Maps API key is restricted to the Android app's
package name and SHA-1 fingerprint.

% ====================================================================
% CHAPTER 3: METHODOLOGY
% ====================================================================
\chapter{Methodology}

This chapter describes the methodology implemented for problem-solving, including the feasibility study, system design, implementation details, and module-level descriptions of every feature in the EdumentX platform.

\section{System Design}

\subsection{System Architecture}

Figure~\ref{fig:sysarc} shows the overall system architecture of
EdumentX. The architecture has three layers: the client
side, the serverless backend, and external services.

\begin{figure}[H]
\centering
\safeincludegraphics[width=\textwidth]{sysarc.png}
\caption{System Architecture of EdumentX}
\label{fig:sysarc}
\end{figure}

The client side is the React Native mobile application
built with Expo. It serves as the single entry point for all
users and communicates with the backend using the native
Firebase SDK.

The serverless backend runs across two platforms. Firebase
provides Authentication (email/password + Google sign-in) and
Cloud Firestore for structured data with Security Rules.
Supabase provides object storage (1 GB free) and Edge Functions
for the RAG chatbot (Groq API), eSewa payment signing, and AI
identity verification (Groq Vision). Firebase Cloud Functions
are scaffolded for Firestore-centric scheduled work.

The external services layer includes Google Maps (via expo-maps
native module on Android), the Groq API for LLM inference
(chatbot and vision), and Nominatim for keyless geocoding.

\subsection{Use Case Diagram}

Figure~\ref{fig:usecase} shows the use case diagram with three
actors: Students/Parents, Tutors, and Admin.

\begin{figure}[H]
\centering
\safeincludegraphics[width=0.75\textwidth]{usecase.jpeg}
\caption{Use Case Diagram}
\label{fig:usecase}
\end{figure}

Students can register, sign in, search for tutors using
filters or the AI chatbot, view tutor profiles, send
enrollment requests, join group batches, message tutors, save
preferred tutors, and rate tutors.

Tutors can register, set their availability, upload
documents for verification, manage their profile, respond to
enrollment requests, create group batches, view their dashboard
with analytics, and upgrade to a Pro subscription.

Admin can log in to the admin panel, review and approve/reject
tutor verifications (with AI pre-screening), manage user
accounts, and monitor platform statistics.

\subsection{Data Flow Diagram}

Figure~\ref{fig:dfd0} shows the Level 0 Data Flow Diagram.

\begin{figure}[H]
\centering
\safeincludegraphics[width=2.8in]{Level_0_DFD.png}
\caption{Level 0 Data Flow Diagram}
\label{fig:dfd0}
\end{figure}

Figure~\ref{fig:dfd1} shows the Level 1 Data Flow Diagram
with detailed system processes including authentication,
tutor search, AI chatbot recommendation, enrollment
management, tutor verification, and rating handling.

\begin{figure}[H]
\centering
\safeincludegraphics[width=\textwidth]{Level_1_DFD.png}
\caption{Level 1 Data Flow Diagram}
\label{fig:dfd1}
\end{figure}

\subsection{RAG Chatbot Workflow}

Figure~\ref{fig:ragworkflow} illustrates the RAG workflow
used by the EdumentX AI chatbot.

\begin{figure}[H]
\centering
\safeincludegraphics[width=\textwidth]{RAG.png}
\caption{RAG-Based AI Chatbot Workflow}
\label{fig:ragworkflow}
\end{figure}

The chatbot works through a six-step pipeline. First, the
student types a natural language query like ``I need a science
tutor near Lalitpur for 3 days a week under Rs 5,000.'' The
full conversation history is included so the chatbot can
handle follow-up questions. Second, the client extracts
structured parameters from the query: subject, location,
budget, days per week, and plan duration. Third, the client
queries the Firestore \texttt{tutors} collection, applies the
bounding-box pre-filter followed by exact Haversine distance,
and applies budget and subject filters. Fourth, the top ten
matching tutor profiles are formatted and injected into the
Groq API prompt as context data. Fifth, the language model
generates a response recommending three to five tutors with
name, rate, distance, rating, and verification status. Sixth,
the response is rendered as a chat message in the app.

\section{Implementation Details}

\subsection{Frontend Development}

The mobile application is built with React Native (v0.81.5)
using TypeScript (v5.9.2) and Expo SDK 54. Styling is handled
through NativeWind 4.2, which maps Tailwind CSS classes to
React Native styles via a custom Babel plugin. Navigation uses
Expo Router v6, a file-based router where each file in
\texttt{src/app/} automatically becomes a route. The design
system uses a custom color palette defined in
\texttt{tailwind.config.js} with tokens such as
\texttt{bg-night}, \texttt{text-amber}, \texttt{bg-surface},
and \texttt{border-border}. Five custom ESLint rules enforce
design token compliance.

\subsection{Backend Development}

The backend is a hybrid architecture. Firebase provides
Authentication and Cloud Firestore with Security Rules.
Supabase provides object storage and three Edge Functions:
\texttt{chat} (RAG chatbot using Groq),
\texttt{create-esewa-order} (eSewa payment signing and
verification), and \texttt{verify-identity} (AI document
verification using Groq Vision). The \texttt{dataSource}
selector pattern allows switching between mock data
repositories and Firebase/Supabase repositories at runtime.

\subsection{Testing Strategy}

Testing is conducted at three levels:

\begin{enumerate}
\item \textbf{Unit tests} (\texttt{npm run test:derived}): 47
tests covering every pure helper function in
\texttt{services/enrollments/derived.ts} and
\texttt{types.ts}.

\item \textbf{Lint rule tests} (\texttt{npm run test:lint-rules}):
21 tests verifying the five custom ESLint rules that enforce
design token compliance.

\item \textbf{Security rules tests} (\texttt{npm run test:rules}):
Six emulator suites exercising Firestore security rules via
the REST API.
\end{enumerate}

\section{Module Implementation}

\subsection{Onboarding and Authentication Module}

The onboarding flow begins with a splash screen displaying the
EdumentX logo, followed by a three-screen carousel introducing
the platform's value propositions. Authentication is handled
by a single screen (\texttt{EmailSignUp.tsx}) that hosts both
sign-up and log-in forms. Email registration calls
\texttt{createUserWithEmailAndPassword} followed by
\texttt{sendEmailVerification}. Google Sign-In uses
\texttt{@react-native-google-signin/google-signin}. The
auth-flow routing in \texttt{\_layout.tsx} runs a 5-step
redirect tree based on authentication state, email
verification, and role assignment.

\subsection{Tutor Profile and Verification Module}

Tutors complete a multi-step profile form. The AI verification
pipeline triggers automatically on profile submission. The
\texttt{verify-identity} Supabase Edge Function performs two
checks using Groq Vision (Qwen 3.6 model): OCR extraction of
the English name from the citizenship certificate, and face
detection on the profile photograph. If both checks pass with
confidence $\geq$ 0.9, the function auto-approves the tutor by
writing to \texttt{tutorVerifications/\{uid\}} via the
Firestore REST API, updating the profile's
\texttt{verificationStatus}, and creating the
\texttt{tutors/\{uid\}} discovery document. Low-confidence
results are routed to manual admin review.

\subsection{Tutor Discovery and Map Module}

The student home screen displays nearby tutors ranked by
Haversine distance with Pro subscription priority. The
\texttt{MapSearch} screen renders Google Maps via
\texttt{expo-maps} with custom teardrop pins showing tutor
avatars. A filter sheet allows students to adjust distance
radius, budget range, subject, grade level, and tutoring mode.
Tutor cards display name, photo, headline, subjects, rating,
monthly rate, distance, and verification badge.

\subsection{Enrollment and Batches Module}

Students send enrollment requests to tutors. Requests are
stored at
\texttt{enrollmentRequests/\{tutorUid\}/requests/\{requestId\}}.
Tutors accept (selecting a time slot), decline, or request more
information. Accepted enrollments are written to
\texttt{enrollments/\{tutorUid\}/roster/\{enrollmentId\}}.
The \texttt{slotKey} format is \texttt{"day:slot"} (e.g.,
\texttt{mon:5-7}). Group batches are managed through a 3-step
wizard with member selection from the live roster.

\subsection{Messaging Module}

In-app messaging uses a 1:1 conversation model. The
\texttt{conversationId} is the sorted participant pair joined
by \texttt{\_\_}, ensuring both parties address the same thread.
Messages include sender, text, timestamp, and read status.
Real-time listeners provide instant delivery with read
receipts.

\subsection{Pro Subscription Module}

Tutors upgrade to Pro through the \texttt{ProUpgradeScreen}.
The eSewa payment flow uses a WebView that auto-submits an
HMAC-SHA256-signed form. Verification performs four checks:
HMAC validation, transaction ledger lookup (replay protection),
amount/product cross-check, and server-to-server status API
confirmation. Pro features include higher capacity ceilings,
more active batches, a distinct Pro badge, a Featured section,
and advanced analytics.

\subsection{Admin Module}

The admin dashboard provides platform statistics using
\texttt{getCountFromServer}. The verification queue displays
pending tutor submissions with document previews, AI
pre-screen results, and approve/reject actions.

% ====================================================================
% CHAPTER 4: SYSTEM DESIGN
% ====================================================================
\chapter{System Design}

\section{Requirement Specification}

\subsection{Functional Requirements}

\textbf{FR1: User Registration and Authentication}
\begin{itemize}
\item Users can register as Student, Tutor, or Admin.
\item Email/password and Google Sign-In authentication.
\item Email verification before dashboard access.
\item Role-based routing to dashboards.
\end{itemize}

\textbf{FR2: Tutor Profile Management}
\begin{itemize}
\item Multi-step profile with subjects, grades, pricing, location.
\item Profile photograph and citizenship document upload.
\item Availability toggle for search visibility.
\end{itemize}

\textbf{FR3: Student Discovery and Filtering}
\begin{itemize}
\item Map-based proximity search with Google Maps.
\item Filters for subject, grade, budget, distance, tutoring mode.
\item Distance-based ranking with Pro subscription priority.
\item Saved tutors functionality.
\end{itemize}

\textbf{FR4: AI-Powered Chatbot}
\begin{itemize}
\item Natural-language queries about available tutors.
\item RAG pipeline using Groq API with Firestore retrieval.
\item Conversation history for follow-up questions.
\end{itemize}

\textbf{FR5: Enrollment Request Handling}
\begin{itemize}
\item Students send requests with preferred time slots.
\item Tutors accept (with slot selection), decline, or request info.
\item Enrollment tracking with start/end dates.
\end{itemize}

\textbf{FR6: Group Batch Management}
\begin{itemize}
\item Tutors create batches with name, subject, capacity, schedule.
\item Member add/remove from live roster.
\item Students browse open batches from enrolled tutors.
\end{itemize}

\textbf{FR7: In-App Messaging}
\begin{itemize}
\item 1:1 conversations with deterministic IDs.
\item Read receipts and real-time delivery.
\end{itemize}

\textbf{FR8: Ratings and Reviews}
\begin{itemize}
\item Star ratings and text reviews for enrolled tutors.
\item Running average computation.
\end{itemize}

\textbf{FR9: AI Identity Verification}
\begin{itemize}
\item Groq Vision OCR for citizenship document name extraction.
\item Groq Vision face detection for profile photograph.
\item Auto-approve at confidence $\geq$ 0.9; manual review otherwise.
\end{itemize}

\textbf{FR10: Pro Subscription}
\begin{itemize}
\item eSewa sandbox payment with HMAC-SHA256 signing.
\item Enhanced capacity, batches, badges, and analytics for Pro.
\end{itemize}

\textbf{FR11: Admin Dashboard}
\begin{itemize}
\item Platform statistics with live aggregation.
\item Verification queue with AI pre-screen results.
\item User management.
\end{itemize}

\subsection{Non-Functional Requirements}
\begin{itemize}
\item \textbf{Usability:} Consistent NativeWind design system with
tactile haptic feedback and accessible color tokens.
\item \textbf{Performance:} Real-time Firestore listeners provide
sub-second updates; map rendering and chatbot responses execute
within acceptable time.
\item \textbf{Security:} Firebase Security Rules enforce owner-only
data access and role-based permissions.
\item \textbf{Privacy:} Location data used solely for proximity
calculations. Student locations masked from non-enrolled tutors.
\item \textbf{Scalability:} Repository pattern with data source
selector allows mock/production switching.
\item \textbf{Maintainability:} Modular architecture with TypeScript
strict mode and custom ESLint rules.
\end{itemize}

\subsection{System Architecture}

EdumentX follows a multi-layer hybrid architecture:

\begin{itemize}
\item \textbf{Presentation Layer:} React Native + NativeWind +
Expo Router mobile application.

\item \textbf{Application Layer:} TypeScript services with repository
pattern. Screen components consume repository interfaces with
runtime data source selection.

\item \textbf{Service Layer:} Firebase Authentication, Cloud
Firestore, Supabase Edge Functions (chat, eSewa, verify-identity),
Supabase Storage.

\item \textbf{Data Layer:} Cloud Firestore for structured data,
Supabase Storage for files, Groq API for AI inference.
\end{itemize}

\subsection{Use Case Diagram}

The principal actors are Student, Tutor, and Admin. Student use
cases include authentication, map-based discovery, AI chatbot,
enrollment requests, batch browsing, messaging, saving tutors,
and reviews. Tutor use cases include profile management, AI
verification, enrollment handling, batch management, dashboard
analytics, and Pro subscription. Admin use cases include
verification review, platform statistics, and user management.

% ====================================================================
% CHAPTER 5: RESULTS AND DISCUSSION
% ====================================================================
\chapter{Results and Discussion}

\section{Software Development Approach}

The project followed an iterative-incremental development model
over approximately six months (February--August 2026). Each phase
produced a working increment that was tested on a real Android
device before proceeding. The mid-term defense (August 2026)
covered the completed foundation: authentication, onboarding,
design system, and admin panel. The post-mid-term phase added
map-based discovery, AI verification, enrollment management,
group batches, messaging, Pro subscription, and the RAG chatbot.

\section{Testing and Evaluation}

\subsection{Unit Test Results}

The enrollment logic test suite (\texttt{npm run test:derived})
covers 47 unit tests across all pure helper functions:

\begin{table}[H]
\centering
\caption{Unit Test Coverage Summary}
\begin{tabular}{|l|c|c|}
\hline
\textbf{Module} & \textbf{Tests} & \textbf{Status} \\
\hline
KTM date helpers & 8 & All passing \\
slotKey parsing & 12 & All passing \\
Today-sessions derivation & 6 & All passing \\
Booked-map generation & 5 & All passing \\
Capacity counting & 8 & All passing \\
Availability draft helpers & 8 & All passing \\
\hline
\textbf{Total} & \textbf{47} & \textbf{All passing} \\
\hline
\end{tabular}
\end{table}

\subsection{Lint Rule Test Results}

The design-token ESLint rule tests (\texttt{npm run test:lint-rules})
verify five custom rules:

\begin{table}[H]
\centering
\caption{ESLint Rule Test Results}
\begin{tabular}{|l|c|c|}
\hline
\textbf{Rule} & \textbf{Tests} & \textbf{Status} \\
\hline
no-raw-hex-placeholder & 4 & All passing \\
no-raw-hex-color-prop & 4 & All passing \\
no-raw-hex-inline-color & 4 & All passing \\
no-non-token-radius & 5 & All passing \\
no-non-token-color-class & 4 & All passing \\
\hline
\textbf{Total} & \textbf{21} & \textbf{All passing} \\
\hline
\end{tabular}
\end{table}

\subsection{Security Rules Test Results}

The Firestore security rules emulator suite (\texttt{npm run test:rules})
exercises six test suites:

\begin{table}[H]
\centering
\caption{Security Rules Test Suites}
\begin{tabular}{|l|c|}
\hline
\textbf{Suite} & \textbf{Status} \\
\hline
Tutors availability carve-out & Passing \\
Batches collections & Passing \\
AcceptRequest transaction paths & Passing \\
Conversations/messages participant gates & Passing \\
Reviews list rules & Passing \\
RemoveEnrollment cascade & Passing \\
\hline
\end{tabular}
\end{table}

\subsection{Build Verification}

The project passes TypeScript type checking
(\texttt{npx tsc --noEmit}) with zero errors and ESLint
(\texttt{npx eslint .}) with zero violations after every
significant change. The Android debug build
(\texttt{npx expo run:android}) completes successfully.

\section{Current Build}

Figure~\ref{fig:mockup1} shows the current build running on an
Android device.

\outputscreen{mockup_1.png}{EdumentX Current Build on Android}{fig:mockup1}

% ====================================================================
% CHAPTER 6: CONCLUSIONS
% ====================================================================
\chapter{Conclusions}

This report presented the design, development, and evaluation of
EdumentX, a location-based mobile application for home tutor
discovery in Kathmandu Valley. The project successfully
addresses the core problems of informal tutor finding, lack of
verification, and absence of proximity-based search through a
comprehensive mobile platform built with React Native and Expo.

The key achievements of this project are:

\begin{itemize}
\item A map-based tutor discovery system using Google Maps with
custom pins and proximity-based ranking, enabling students to
find nearby tutors within the Kathmandu Valley.

\item An AI-powered RAG chatbot using Groq that answers
natural-language queries about available tutors, grounded in
live Firestore data.

\item An automated identity verification pipeline using Groq
Vision OCR for citizenship document validation and face
detection for profile photographs, with automatic approval for
high-confidence matches and manual admin review for ambiguous
cases.

\item A structured enrollment management system with slot-based
scheduling, enrollment request workflows, and group batch
management.

\item An in-app messaging system with real-time delivery and
read receipts.

\item A Pro subscription model with eSewa payment integration
(sandbox), providing enhanced visibility and features for
subscribing tutors.

\item A complete admin dashboard with verification queues,
platform statistics, and user management.

\item A zero-cost architecture operating within Firebase Blaze
and Supabase free-tier quotas, making the platform viable as
a demonstration prototype without paid infrastructure.
\end{itemize}

The project demonstrates that a location-first, trust-oriented
tutoring marketplace can be built and deployed using modern
cross-platform frameworks and serverless services at zero
operational cost at demo scale.

% ====================================================================
% CHAPTER 7: LIMITATIONS AND FUTURE ENHANCEMENT
% ====================================================================
\chapter{Limitations and Future Enhancement}

\section{Limitations}

\begin{enumerate}
\item \textbf{Android-only:} The current build targets Android
exclusively. iOS support requires an Apple Developer account
and additional testing with Apple Maps integration via
expo-maps.

\item \textbf{eSewa sandbox:} The payment integration uses
eSewa's sandbox environment with test credentials
(\texttt{EPAYTEST}). Production deployment requires a
registered merchant account and live credentials.

\item \textbf{No video tutoring:} The platform supports only
in-person home tutoring. Online or video-based classes are
outside the current scope.

\item \textbf{Limited tutor coverage:} The platform currently
has a small number of registered tutors. A critical mass of
tutors is needed before the recommendation engine and chatbot
provide meaningful results.

\item \textbf{HuggingFace DNS:} The HuggingFace Serverless
Inference API is unreachable from Supabase Edge Runtime due
to DNS limitations. Face detection uses Groq Vision as a
fallback.

\item \textbf{No push notifications:} While expo-notifications
is installed, FCM push notifications for real-time alerts
(messaging, enrollment updates) are not yet fully implemented.
The current build uses local notifications only.
\end{enumerate}

\section{Future Enhancement}

\begin{enumerate}
\item \textbf{iOS support:} Extend the build to iOS with Apple
Maps integration and test on physical iOS devices.

\item \textbf{FCM push notifications:} Implement server-side
push notifications via Firebase Cloud Functions for messaging,
enrollment updates, and verification status changes.

\item \textbf{Online tutoring:} Add video call integration
(e.g., WebRTC) for online tutoring sessions.

\item \textbf{Tutor analytics:} Expand the analytics dashboard
with profile view trends, subject demand patterns, and
conversion rates.

\item \textbf{Multi-language support:} Add Nepali language
support for the UI and chatbot responses.

\item \textbf{EAS Build and distribution:} Set up EAS Build
for automated Android builds and distribute via Firebase App
Distribution for testing.
\end{enumerate}

% ====================================================================
% REFERENCES
% ====================================================================
\chapter{References}
\bibliographystyle{unsrt}
\bibliography{ref}
\end{document}
